\documentclass[11pt]{article}

\usepackage[margin=1in]{geometry}
\usepackage{amsmath, amssymb}
\usepackage{enumitem}
\usepackage{parskip}
\usepackage[colorlinks=true, linkcolor=blue, urlcolor=blue]{hyperref}
\usepackage{xcolor}

% lightweight environment for the meta/marginal notes 
\newcommand{\metanote}[1]{\par\noindent\textcolor{gray}{\textit{#1}}\par}

\setlist{itemsep=2pt, topsep=4pt}

\title{MATS week 1 proposal}
\author{Alessa Carbo}
\date{\today}

\begin{document}

\maketitle



\section*{Problem Statement}

Current machine unlearning approaches rely on behavioral output evaluations or
proxy-based measurements. The field keeps shipping unlearning/tamper-resistance
methods that report near-complete suppression and then get red-teamed for
dollars of compute. Cas documents that many
``tamper-resistant'' losses can be gamed by simple rotation, contraction, or
translation of activations. In \emph{Deep Ignorance} (O'Brien, Casper et al.,
2025) it is shown that filtering bio content from \emph{pretraining} data
resists ${\sim}10\text{k}$ steps / 300M tokens of adversarial fine-tuning, which
is an order of magnitude more --- but the same models recover the capability
when the information is supplied in context (retrieval/search), so even the
strongest weight-space defense has an axis on which it fails.

What appears to be missing is an actual measurement tool that distinguishes
``destroyed'' from ``hidden behind a cheap transform,'' ideally without having
to spend the full adversarial budget, and doing so by looking \emph{inside} the
model rather than solely relying on its outputs.

\section*{Contribution}

To build a set of measurement instruments capable of quantifying how much of a
suppressed capability is actually \emph{destroyed} in an unlearned model vs.\
merely hidden by a cheap transformation, and to use this set of instruments to
evaluate the current zoo of unlearning algorithms on bio-risk-relevant
capability suppression.

\subsection*{Sketch / attempt at formalization of what this could look like}

\[
  \mathrm{Hid}_\tau(M, T) \;=\; \min\bigl\{\, c \;:\; \mathrm{Acc}\!\left(R_c(M),\, T\right) \ge \tau \,\bigr\}
\]

Where:
\begin{itemize}
  \item $M$: unlearned model
  \item $T$: target task
  \item $R$: recovery procedure
  \item $c$: budget
  \item $\tau$: threshold
\end{itemize}

The idea would be that:
\begin{itemize}
  \item \textbf{Low Hid}: corresponds to shallowness (maybe a linear probe / few steps suffice).
  \item \textbf{High Hid}: deep (could need large capacity, many steps, or never crosses).
\end{itemize}

Here we could have a deliverable which shows a recovery curve --- recovered
capability vs.\ budget, per model --- plus this scalar summary.

Some important things to note: this formalization notion and the quantities
composing it are quite flexible. For example, our budget $c$ is not one thing.
We can think of at least 3 candidates now:

\begin{enumerate}
  \item \textbf{Representational / probe capacity}: at frozen weights, what
  computational class can decode the answer from the activations? (e.g., what
  kind of probe --- linear, MLP depth, attention-pooled).
  \item \textbf{Optimization / weight change}: how many fine-tuning steps (or
  weight movement, however we want to define it) does it take to reassemble the
  capability?
  \item \textbf{In-context budget}: how many ICL demonstrations does it take to
  reconstitute it? We already know even filtering can fail at this.
\end{enumerate}

It'd be interesting to see what type of relationship we have between axes, like
whether there is a separation --- such as a capability being hard to recover on
the weight axis while easy on the ICL one. It's also still unclear whether
``recovery-complexity'' has genuine class structure or it's simply a continuous
``axis.''

\section*{Proposed Methodology / Experiments}

\subsection*{Experiment 1: Probing --- zero-weight-change recovery}

\textbf{Battery of models:} RMU, RepNoise, Circuit Breakers, Deep Ignorance
filtered models.

\textbf{Data:} WMDP forget set; probes trained on residual-stream activations.
\begin{itemize}
  \item Linear probes per layer
  \item MLP probes at different depths
  \item Attention-pooled probes (in case we want to attend to a specific token position)
\end{itemize}

\textbf{What this experiment could show:} a per-model, per-layer, per-capacity
profile of recoverable WMDP-bio info that would be plotted / compared to WMDP
accuracy. The intuition is something like: high probe accuracy on layers where
behavioral accuracy is floored $\rightarrow$ info is present, and suppression is
shallow. Filtered models could serve as a type of control: if filtering truly
removed knowledge, we would expect probe accuracy to stay low even at high
capacity. Here capacity would refer to how much computation the probe is allowed
to do when reading the answer off the frozen activations. A linear probe can
only apply one affine map; an MLP probe with hidden layers can do more. A
higher-capacity probe would be one that can implement a richer class of
functions from activations to answer. (Note this can change, and maybe it's not
that important right now for us to have a super-precise definition of high
capacity, but simply to deploy more than one type of probe.)

\textbf{Control variables:}
\begin{enumerate}
  \item \emph{Model where knowledge was never present.} We would want to compare
  each probe against the same architecture trained on a model that never had the
  knowledge. Because if our probe also scores high on this control, then it might
  be learning the task from its own labels and not reading from hidden knowledge.
  This is a contamination risk factor to keep in mind.
  \item \emph{Selectivity.} We'd report the probe profile on the forget dataset,
  a retain set, and some unknown control task.
\end{enumerate}

\textbf{Causal steering:} once we have probes that can locate whether suppressed
info is linearly present, we'd want to test whether adding the probe direction
back into the residual stream recovers a capability at zero fine-tuning.

\subsection*{Experiment 2: Recovery dynamics}

\textbf{Battery of models:} same as Experiment 1.

Now we would run unlearning and adversarial fine-tuning on held-out WMDP-style
data. Per step we would record:
\begin{itemize}
  \item Behavioral WMDP accuracy
  \item Probe accuracies from our first experiment
  \item Weight changes
  \item Loss curve
  \item Representation similarity (CKA, Procrustes, $R^2$, or some metric
  inspired by the model-stitching literature)
  \item Activation-rank difference across the forget set
\end{itemize}

\textbf{What we would want to show with this experiment:} does there exist an
early cheap-and-reversible phase, then a later durable phase? The intuition would
be something like distinguishing a wrapper-revival vs.\ genuine reconstruction
(representations move substantially).

\textbf{Self-bootstrapping recovery:} I need a little more time to develop
this\ldots\ but I have some fuzzy ideas.

\subsection*{Experiment 3: Mutual Information instrument}

Adapting from Nakkiran's conditional-MI construct to the unlearning setting:

\[
  Q(M, T, b) \;=\; I(M_T;\, Y) \;-\; I\!\left(M_T;\, Y \mid R_b(M)_T\right)
\]

Where:
\begin{itemize}
  \item $Q$ = the part of $M$'s residual capability which is \emph{not}
  explained by the budget-$b$ recovery.
\end{itemize}

With $Q \approx 0$: whatever's left, the cheap recovery already finds it
(shallow). $Q$ large: $M$ holds capability the cheap recovery can't reach ---
maybe suppression worked, or the recovery is too weak.

\textbf{Calibration:} does low $Q$ at budget $b$ predict that a real attacker at
budget $b$ recovers the capability? Ideally we would want this predictive
capability, because if not then there's not much utility or added value in this
measurement.

\textbf{Note on data regime:} for binary MCQ correctness, $I(M_T; Y)$ is
monotone with accuracy, so $Q$'s only content beyond a plain accuracy-gap is the
\emph{conditioning}. In many cases $Q = (\text{full-recovery accuracy}) -
(\text{budget-}b\text{ recovery accuracy})$, a curve gap. Run Exp 3 specifically
to find whether $Q$ ever disagrees with the accuracy gap, it can, when $M$
and $R_b(M)$ make correlated vs.\ independent errors, since MI sees the joint
distribution. If a regime exists where two models have equal recovered accuracy
but different $Q$ because error structure differs, MI seems very noteworthy. If
not, demote it to ``we also report an information-theoretic version that agrees
with the gap'' and lead with the gap.

\subsection*{Connecting to Cas'  doc and stated hypotheses for what method is already doing/failing at }

A useful way to read these experiments is as adjudicating specific failure modes
that Cas (2026) has already \emph{hypothesized} for several of these methods.
, several tamper-resistant losses are suspected of reducing
to a cheap transformation that does not genuinely destroy information and these
are hypotheses stated from the loss-function design, and the instruments here
are what would confirm or refute them mechanistically, per method.

\textbf{Circuit Breakers --- hypothesized conditional orthogonal rotation.}
Cas notes that CB trains an adapter to make forget-set activations orthogonal
to their original directions, and suspects this can be satisfied by learning
``a simple conditional orthogonal rotation'' that should be ``removed easily by
further fine-tuning.'' If that hypothesis holds, the suppressed direction still
exists and is merely rotated: our causal-steering test should recover capability
at zero fine-tuning by adding the probe direction back, and our
representation-similarity measures during recovery (Exp 2) should show a small,
structured, low-rank change with high similarity to the unlearned model
(undoing a rotation), rather than substantial representational movement.

\textbf{RepNoise --- hypothesized contraction, not destruction.}
Cas notes that RepNoise trains forget activations toward white noise via MMD,
and objects that ``having a Gaussian shape does not imply being genuinely
random'' --- the procedure ``will probably just contract the distribution of
latents'' rather than scrub information. If that hypothesis holds, the
information survives the contraction: our probe-capacity profile (Exp 1) should
still decode the answer at sufficient capacity even where behavioral accuracy is
floored, and our activation-rank measures (Exp 2) should carry a contraction
signature distinct from the rotation signature above. Note this is a
\emph{different} suspected failure mode than CB's --- so it is also a test of
whether the instrument discriminates between mechanisms.

\textbf{The scaling prediction --- ``cheap solution vs.\ genuinely hard
objective.''} Cas predicts that scaling an unlearning algorithm (running it
much longer) only helps if the forget objective ``is genuinely hard to learn and
cannot be minimized with a simple solution,'' and conjectures that a
characteristic of good unlearning is that it \emph{benefits from scaling}. Our
recovery-dynamics experiment operationalizes exactly this distinction: the
early-cheap-reversible vs.\ later-durable phase structure is a mechanistic proxy
for ``cheap solution vs.\ genuinely hard objective.'' This means Exp 2 could
predict which methods will benefit from his longer-training idea \emph{before}
committing the compute to run it.

\textbf{Methods without a pre-stated mechanism (e.g.\ RMU).} Not every method in
the battery has a specific failure-mode hypothesis yet and
the instrument is exploratory, characterizing what the method does internally
ather than adjudicating a stated claim. This is worth keeping explicit so we
don't retrofit a hypothesis where none was stated.

\section*{4. Characterizing Meta-Unlearning}

\metanote{Thanks Eric for helping sketch out this more concretely!}

For protected distribution $D^*$, maximize loss on $D^*$ not only at the current
model but across all future models reachable by adaptation, where each future
model is produced by an inner argmin on data the current model collects.

\[
  \max \;\;
  \underbrace{\mathbb{E}_{(x,y)\sim D^*}\!\left[\ell(\theta_t; x, y)\right]}_{\text{loss now}}
  \;+\;
  \underbrace{\sum_j \mathbb{E}_{(x,y)\sim D^*}\!\left[\ell(\theta_{t+j}; x, y)\right]}_{\text{loss at all future models}}
  \qquad \text{s.t.} \qquad
  \theta_{t+1} = \arg\min_{\theta}\; \mathbb{E}_{(x,y)\sim D(\theta_t)}\!\left[\ell(\theta; x, y)\right]
\]

Where:
\begin{itemize}
  \item $D^*$: protected data distribution --- the thing you want the model to
  stay bad at. Distribution over (input, label) pairs for the dangerous
  capability. This is what we want high loss on.
  \item $(x, y)$: an input and its label drawn from $D^*$.
  \item $\theta_t$: the model parameters at round $t$. $\theta_t$ is the output
  as defender --- the unlearned model you're trying to choose. This is the
  variable the outer max optimizes.
  \item $\ell(\theta; x, y)$: the per-example loss of the model with parameters
  $\theta$ on example $(x,y)$. High $\ell$ on $D^*$ = the model being bad at the
  protected task = good from the defender's point of view. This is why we have a
  max here, not a min.
  \item $\theta_{t+j}$: the model parameters $j$ rounds into the future, produced
  by the recovery recursion below.
  \item $\sum_j$: sum over future rounds. The defender wants the model to stay
  bad at $D^*$ not just now but at every future model the adversary can reach.
  \item $D(\theta_t)$: the distribution map --- the data distribution that gets
  collected/generated when the model is $\theta_t$. The key feature is that it's
  a function of the current model, as the model influences the data its successor
  trains on.
  \item $S$ (ICL variant): a set of in-context demonstrations. In the ICL
  version, the adversary doesn't change weights; the inner step searches over
  demonstration sets $S$ that, placed in context, recover capability.
\end{itemize}

\textbf{Inner recursion:}
\[
  \theta_{t+1} = \arg\min_{\theta}\; \mathbb{E}_{(x,y)\sim D(\theta_t)}\!\left[\ell(\theta; x, y)\right]
\]
This is the adversary's move: given the current model $\theta_t$, the adversary
collects data from $D(\theta_t)$ and retrains to minimize loss on it, producing
the next model $\theta_{t+1}$. This is standard learning (yes, minimizing loss)
pointed at recovering the capability. The defender is trying to pick $\theta_t$
so that even after the adversary runs this, loss on $D^*$ stays high.

\end{document}